\documentclass{article}
\usepackage{amsmath,amssymb}
\usepackage[vocab={graphcounts}]{leanify}
\newtheorem{lemma}{Lemma}
\begin{document}

\section{Notation}

We write $e(n)$ for the number of edges of the complete graph on $n$ vertices.
We write $b(m,n)$ for the number of edges of the complete bipartite graph with parts of sizes $m$ and $n$.
We write $q(n)$ for the number of edges of the hypercube of dimension $n$.
We write $L(n)$ for the number of labelled graphs on $n$ vertices.
We write $\tau(n)$ for the number of spanning trees of the complete graph on $n$ vertices.
We write $\sigma(m,n)$ for the number of spanning trees of the complete bipartite graph with parts of sizes $m$ and $n$.
We write $c(n)$ for the number of spanning trees of the cycle on $n$ vertices.
We write $P(n,k)$ for the chromatic polynomial of the complete graph on $n$ vertices at $k$.
We write $T(n,k)$ for the chromatic polynomial of a tree on $n$ vertices at $k$.
We write $C(n,k)$ for the chromatic polynomial of the cycle on $n$ vertices at $k$.
We write $t(n,r)$ for the number of edges of the Tur\'an graph on $n$ vertices with $r$ parts.
We write $R(a,b)$ for the binomial bound on the Ramsey number for $a$ and $b$.

\section{Counts}

\begin{lemma}\label{edges}
$e(4)=6$, $b(2,3)=6$, $q(3)=12$ and $L(3)=8$.
\end{lemma}

\begin{lemma}\label{trees}
$\tau(4)=16$, $\sigma(2,2)=4$ and $c(5)=5$.
\end{lemma}

\section{Colourings and extremal numbers}

\begin{lemma}\label{chromatic}
$P(3,3)=6$, $T(3,3)=12$, $C(3,3)=6$ and $C(4,2)=2$.
\end{lemma}

\begin{lemma}\label{factorial}
Let $n$ be an integer with $n\ge2$. Then $P(n,n)=n!$.
\end{lemma}

\begin{lemma}\label{turan}
$t(5,2)=6$, $t(6,3)=12$ and $R(3,3)=6$.
\end{lemma}

\section{Closed forms}

\begin{lemma}\label{cayley}
For every positive integer $n$, $\tau(n)=n^{n-2}$. For every nonnegative integer $n$, $e(n)=\frac{n(n-1)}{2}$ and $L(n)=2^{e(n)}$.
\end{lemma}

\begin{lemma}\label{trees-chromatic}
Let $n$ be a positive integer and let $k$ be a real number. Then $T(n,k)=k(k-1)^{n-1}$ and $P(n,k)=\prod_{i=0}^{n-1}(k-i)$.
\end{lemma}

\end{document}
